\section{The parameter inequalities}\label{cc:app:parameters}
\begin{proof}[Proof of Lemma~\ref{cc:uf:tower:lem:slow-orders}]
The current order $q_j$ is always admissible in the maximum
\eqref{cc:uf:tower:eq:schedule}. Indeed, $\mathcal B(q_1)=\mathcal B(1)
\le N=\lambda_0$, and for $j\ge2$ the preceding choice gives
$\mathcal B(q_j)\le\lambda_{j-2}\le\lambda_{j-1}$. Thus $q_{j+1}=q_j$
until $\lambda_{j-1}$ reaches $\mathcal B(q_j+1)$, at which point the
order increases by one. Since $N>1$ and $K_j\ge2$, the frequencies tend
to infinity. They eventually exceed every fixed threshold, so
$q_j\to\infty$.

The gain $L_{j+1}$ involves three successive frequency steps. By
\eqref{cc:uf:tower:eq:schedule},
\[
 L_{j+1}=4\log\lambda_{j+2}
        =4K_jK_{j+1}K_{j+2}\log\lambda_{j-1}.
\]
Each exponent is at most $K(q_{j+1}+1)$, since the orders increase
by at most one. Thus \eqref{cc:uf:tower:eq:uniform} follows from
\eqref{cc:par:floor} and the choice of $q_{j+1}$, which give
\[
 4K_jK_{j+1}K_{j+2}\le4K(q_{j+1}+1)^3
 \le\log\mathcal B(q_{j+1})\le\log\lambda_{j-1}. \qedhere
\]
\end{proof}
